\documentclass{article}
\linespread{1.3}
\usepackage[margin=50pt]{geometry}
\usepackage{amsmath, amsthm, amssymb, amsthm, tikz, fancyhdr, float, graphicx, spverbatim, systeme}
\pagestyle{fancy}
\renewcommand{\headrulewidth}{0pt}
\newcommand{\changefont}{\fontsize{15}{15}\selectfont}

\fancypagestyle{firstpageheader}{
  \fancyhead[R]{
    \changefont
    \parbox[t]{4cm}{ % Adjust width as needed
      Michael Huang\\
      EN.555.644.81\\
      Homework 3
    }
  }
}

\begin{document}

\thispagestyle{firstpageheader}
{\Large 

\section*{4.5}

Suppose that zero interest rates with continuous compounding are as follows:
\begin{center}
\begin{tabular}{ |c|c| } 
 \hline
 Maturity (months) & Rate (\% per annum) \\ 
 \hline
 3 & 8.0 \\ 
 \hline
 6 & 8.2 \\ 
 \hline
 9 & 8.4 \\ 
 \hline
 12 & 8.5 \\ 
 \hline
 15 & 8.6 \\ 
 \hline
 18 & 8.7 \\ 
 \hline
\end{tabular}
\end{center}
Calculate forward interest rates for the second, third, fourth, fifth, and sixth quarters. \\ \\
We know that the forward rate for the period between \(T_1\) and \(T_2\) implied by today's curve is \(R_F = \frac{R_2 T_2 - R_1 T_1}{T_2 - T_1}\), so we calculate \(R_F\) accordingly, with each quarter being \(0.25T\).

\begin{center}
\begin{tabular}{ |c|c| } 
 \hline
 Quarter & Forward interest rate \\ 
 \hline
 2 & \(\frac{0.082(0.5) - 0.08(0.25)}{0.25} = \frac{0.041 - 0.02}{0.25} = \frac{0.021}{0.25} = 0.084\), or \textbf{8.4\%} \\ 
 \hline
 3 & \(\frac{0.084(0.75) - 0.082(0.5)}{0.25} = \frac{0.063 - 0.041}{0.25} = \frac{0.022}{0.25} = 0.088\), or \textbf{8.8\%} \\ 
 \hline
 4 & \(\frac{0.085(1.00) - 0.084(0.75)}{0.25} = \frac{0.085 - 0.063}{0.25} = \frac{0.022}{0.25} =  0.088\), or \textbf{8.8\%} \\ 
 \hline
 5 & \(\frac{0.086(1.25) - 0.085(1.00)}{0.25} = \frac{0.1075 - 0.085}{0.25} = \frac{0.0225}{0.25} = 0.09\), or \textbf{9.0\%} \\ 
 \hline
 6 & \(\frac{0.087(1.50) - 0.086(1.25)}{0.25} = \frac{0.1305 - 0.1075}{0.25} = \frac{0.023}{0.25} = 0.092\), or \textbf{9.2\%} \\ 
\hline
\end{tabular}
\end{center}

}
{\Large 

\section*{4.8}

What does duration tell you about the sensitivity of a bond portfolio to interest rates? What are the limitations of the duration measure? \\ \\ 
Duration measures the \textbf{bond portfolio's value sensitivity relative to to interest rate changes on the yield curve}. By definition, we use the duration to calculate the amount that the value of the portfolio changes: the percentage decrease in the value of the bond portfolio equals the duration multiplied by the size of the shift on the yield curve. The duration is limited as a measure in that the shift on the yield curve must \textbf{both small and parallel}, with all maturities moving by the same amount. 

}
{\Large 

\section*{4.9}

What rate of interest with continuous compounding is equivalent to 15\% per annum with monthly compounding? \\ \\ 
We solve using the corresponding equation for continuous compounding, with \(R_m = 0.15, m = 12\):
\[
  e^{R_c} = (1 + \frac{R_m}{m})^m
\]
\[
  e^{R_c} = (1 + \frac{0.15}{12})^{12}
\]
\[
  e^{R_c} = 1.0125^{12}
\]
\[
  R_c = \ln(1.0125^{12})
\]
\[
  R_c = 12 \ln(1.0125)
\]
\[
  R_c \approx 0.14907023988
\]
which indicates that the rate is about \boxed{\textbf{14.91\%}}.


}
{\Large 

\section*{4.11}

Suppose that 6-month, 12-month, 18-month, 24-month, and 30-month zero rates are 4\%, 4.2\%, 4.4\%, 4.6\%, and 4.8\% per annum with continuous compounding respectively. Estimate the cash price of a bond with a face value of 100 that will mature in 30 months and pays a coupon of 4\% per annum semiannually. \\ \\
We need to discount each cash flow at the matching zero rate. We know that the coupon pays \(\frac{0.04 \cdot 100}{2} = \) = \$2 every 6 months, and then 100 + 2 = \$102 at the last payment segment. We therefore would calculate the price like so:
\[
  P = 2e^{-0.04(0.5)} + 2e^{-0.042(1.0)} + 2e^{-0.044(1.5)} + 2e^{-0.046(2)} + 102e^{-0.048(2.5)}
\]
\[
  = 1.96039734661 + 1.91773956114 + 1.87226172858 + 1.82421029909 + 90.4658845452
\]
\[
  P = 98.0404934806
\]
So the price is approximately \boxed{\textbf{\$98.04}}.

}
{\Large 

\section*{4.12}

A three-year bond provides a coupon of 8\% semiannually and has a cash price of 104. What is the bond's yield? \\ \\
We know that the coupon pays \(\frac{0.08 \cdot 100}{2} = \) \$4 due to the semiannual payments. Given that the bond itself therefore pays \$4 every 6 months until maturity, at which it pays \$104, we need to therefore solve 
\[
  4e^{-0.5y} + 4e^{-y} + 4e^{-1.5y} + 4e^{-2y} + 4e^{-2.5y} + 104e^{-3y} = 104
\]
The textbook solves this using trial and error, which I would prefer not to, but let's give it a try:
\begin{center}
\begin{tabular}{ |c|c| } 
 \hline
 Guess for \(y\) & Computed price \\ 
 \hline
 0.06 & 105.163180649 \\ 
 \hline
 0.065 & 103.736019541 \\ 
 \hline
 0.0625 & 104.447023845 \\ 
 \hline
 0.0635 & 104.162005945 \\ 
 \hline
 0.064 & 104.019805386 \\ 
 \hline
 0.06425 & 103.948782073 \\ 
 \hline
 0.0641 & 103.991389907 \\
 \hline
 0.06405 & 104.005596621 \\
 \hline
\end{tabular}
\end{center}
We're extremely close at this point. It seems reasonable enough to say that the rate appears to be very close to 0.06405, or \boxed{\textbf{6.405\%}}.

}
{\Large 

\section*{4.14}

Suppose that zero interest rates with continuous compounding are as follows:
\begin{center}
\begin{tabular}{ |c|c| } 
 \hline
 Maturity (years) & Rate (\% per annum) \\ 
 \hline
 1 & 2.0 \\ 
 \hline
 2 & 3.0 \\ 
 \hline
 3 & 3.7 \\ 
 \hline
 4 & 4.2 \\ 
 \hline
 5 & 4.5 \\ 
 \hline
\end{tabular}
\end{center}
Calculate forward interest rates for the second, third, fourth, and fifth years. \\ \\
We know that the forward rate for the period between \(T_1\) and \(T_2\) implied by today's curve is \(R_F = \frac{R_2 T_2 - R_1 T_1}{T_2 - T_1}\). We are always going year to year in this problem, with each \(T\) being a year, so we can further simplify this to just \(R_F = R_2T_2 - R_1T_1\).

\begin{center}
\begin{tabular}{ |c|c| } 
 \hline
 Year & Forward interest rate \\ 
 \hline
 2 & \(0.03(2) - 0.02(1) = 0.06 - 0.02 = 0.04\), or \textbf{4.0\%} \\ 
 \hline
 3 & \(0.037(3) - 0.03(2) = 0.111 - 0.06 = 0.051\), or \textbf{5.1\%} \\ 
 \hline
 4 & \(0.042(4) - 0.037(3) = 0.168 - 0.111 = 0.057\), or \textbf{5.7\%} \\ 
 \hline
 5 & \(0.045(5) - 0.042(4) = 0.225 - 0.168 = 0.057\), or \textbf{5.7\%} \\ 
 \hline
\end{tabular}
\end{center}

}
{\Large 

\section*{4.16}

A 10-year, 8\% coupon bond currently sells for \$90. A 10-year, 4\% coupon bond currently sells for \$80. What is the 10-year zero rate? (Hint: Consider taking a long position in two of
the 4\% coupon bonds and a short position in one of the 8\% coupon bonds.) \\ \\
Using the hint, we long two of the 4\% bonds and short one of the 8\% bonds. These balance out, with the coupon payment from the longs being \(\frac{0.04 \cdot 100}{1} = \) \$4 each, so we get \(2 \cdot \) \$4 = \$8. For the short, we pay \(\frac{0.08 \cdot 100}{1} = \) \$8. So the values amounts cancel out to 0 dollars each year from coupons. So looking at the actual cash flows, we only consider the initial bond transactions and the final expiry date. For the initial transaction: 
\[
  2 \cdot \text{-\$80} + 1 \cdot \text{\$90} = \text{-\$160 + \$90 = -\$70}
\]
And for the expiry date:
\[
  2 \cdot \text{\$100} - 1 \cdot \text{\$100} = \text{\$200 - \$100 = \$100}
\]
We treat this as a synthetic zero-coupon as the intermediate cash flows are 0, so we can solve the rate as 
\[
  70 \cdot e^{10(R_c)} = 100
\]
\[
  e^{10(R_c)} = \frac{10}{7}
\] 
\[
  10(R_c) = \ln(\frac{10}{7})
\]
\[
  R_c = \frac{\ln(\frac{10}{7})}{10} = \frac{0.35667494393}{10} = 0.035667494393
\]
So in conclusion the overall rate is approximately \boxed{\textbf{3.57\%}}.

}
{\Large 

\section*{4.22}

A five-year bond with a yield of 11\% (continuously compounded) pays an 8\% coupon at the end of each year.

\subsection*{a)}

What is the bond's price? \\ \\


\subsection*{b)}

What is the bond's duration? \\ \\


\subsection*{c)}

Use the duration to calculate the effect on the bond’s price of a 0.2\% decrease in its yield. \\ \\


\subsection*{d)}

Recalculate the bond's price on the basis of a 10.8\% per annum yield and verify that the result is in agreement with your answer to (c). \\ \\ 


}
{\Large 

\section*{4.34}

The following table gives the prices of bonds
\begin{center}
\begin{tabular}{ |c|c|c|c| } 
 \hline
 Bond Principal (\$) & Time to Maturity (yrs) & Annual Coupon (\$)* & Bond Price (\$) \\ 
 \hline
 100 & 0.5 & 0.0 & 98 \\ 
 \hline
 100 & 1.0 & 0.0 & 95 \\ 
 \hline
 100 & 1.5 & 6.2 & 101 \\ 
 \hline
 100 & 2.0 & 8.0 & 104 \\ 
 \hline
\end{tabular}
\end{center}
*Half the stated coupon is paid every six months.

\subsection*{a)}
Calculate zero rates for maturities of 6 months, 12 months, 18 months, and 24 months. \\ \\


\subsection*{b)}
What are the forward rates for the periods: 6 months to 12 months, 12 months to 18 months, 18 months to 24 months? \\ \\ 


\subsection*{c)}
What are the 6-month, 12-month, 18-month, and 24-month par yields for bonds that provide semiannual coupon payments? \\ \\


\subsection*{d)}
Estimate the price and yield of a two-year bond providing a semiannual coupon of 7\% per annum. \\ \\


}

\end{document}